\documentclass[11pt]{article}

% Packages
\usepackage{amsmath}
\usepackage{amssymb}
\usepackage{amsthm}
\usepackage{geometry}
\usepackage{hyperref}
\usepackage{listings}
\usepackage{xcolor}
\usepackage{graphicx}
\usepackage{booktabs}
\usepackage{enumitem}
\usepackage{mathtools}

% Page geometry
\geometry{a5paper,margin=2cm}

% Hyperref setup
\hypersetup{
    colorlinks=true,
    linkcolor=blue!60!black,
    citecolor=green!50!black,
    urlcolor=blue!70!black,
}

% Code listing style
\lstdefinestyle{python}{
    language=Python,
    basicstyle=\ttfamily\small,
    keywordstyle=\color{blue},
    stringstyle=\color{red!60!black},
    commentstyle=\color{green!50!black}\itshape,
    numbers=left,
    numberstyle=\tiny\color{gray},
    stepnumber=1,
    numbersep=8pt,
    frame=single,
    backgroundcolor=\color{gray!8},
    breaklines=true,
    showstringspaces=false,
    tabsize=4,
}

% Theorem environments
\newtheorem{theorem}{Theorem}[section]
\newtheorem{definition}[theorem]{Definition}
\newtheorem{remark}[theorem]{Remark}

% Custom commands
\newcommand{\R}{\mathbb{R}}
\newcommand{\dx}{\,\mathrm{d}x}
\newcommand{\ds}{\,\mathrm{d}s}
\newcommand{\dS}{\,\mathrm{d}S}
\newcommand{\dt}{\,\mathrm{d}t}
\newcommand{\Div}{\nabla\cdot}
\newcommand{\Grad}{\nabla}
\newcommand{\Lapl}{\Delta}
\newcommand{\inner}[2]{\langle #1,\, #2 \rangle}
\newcommand{\norm}[1]{\lVert #1 \rVert}
\newcommand{\domain}{\Omega}
\newcommand{\boundary}{\partial\Omega}

% ---------------------------------------------------------------
\title{%
  \textbf{FEniCSx: Theory and Practice}\\[0.4em]
  \large A Didactic Introduction to the Finite Element Method\\
  and Its Implementation in FEniCSx
}
\author{Will Robertson, with the help of AI}
\date{\today}
% ---------------------------------------------------------------

\begin{document}

\maketitle
\begin{abstract}
\em
This document is intended as a truly `soft ramp' into the finite element method
using FEniCSx.
The current version of the document has been largely written using AI but I (WR)
intend to expand and supplement and modify this material over time\dots
\end{abstract}
\tableofcontents
\newpage

% ================================================================
\section{Introduction}
\label{sec:introduction}
% ================================================================

The \textbf{Finite Element Method} (FEM) is one of the most powerful and
widely used numerical techniques for solving partial differential equations
(PDEs) arising in science and engineering.
\textbf{FEniCSx} is a modern, open-source computing platform for solving PDEs
using the finite element method~\cite{LoggMardalWells2012, Scroggs2022}.
It provides high-level abstractions that allow users to express variational
formulations close to their mathematical description, while transparently
generating efficient low-level code.

\subsection{What is FEniCSx?}

FEniCSx is the latest generation of the FEniCS Project, succeeding the
original FEniCS library.
The project consists of several components:
\begin{itemize}
  \item \textbf{DOLFINx} -- the core C++/Python library providing data
    structures for meshes, function spaces, and linear algebra.
  \item \textbf{UFL} (Unified Form Language) -- a domain-specific language
    for expressing variational forms in near-mathematical notation.
  \item \textbf{FFCx} (FEniCS Form Compiler X) -- compiles UFL forms to
    efficient C code.
  \item \textbf{Basix} -- a finite element tabulation library providing
    basis functions for a wide range of element families.
\end{itemize}

\subsection{Scope of This Document}

This document provides a didactic introduction to the mathematical foundations
underlying FEniCSx and their mapping to FEniCSx concepts and terminology.
It is structured as follows:
\begin{enumerate}
  \item Section~\ref{sec:math} reviews the mathematical foundations:
    function spaces, variational formulations, and the Galerkin method.
  \item Section~\ref{sec:implementation} describes how these mathematical
    constructs map to FEniCSx data structures and workflow.
  \item Section~\ref{sec:examples} presents worked examples that can be
    run directly with FEniCSx.
\end{enumerate}

% ================================================================
\section{Mathematical Foundations}
\label{sec:math}
% ================================================================

\subsection{What Types of Problems Are Represented by PDEs?}
\label{subsec:pde-types}

Partial differential equations (PDEs) arise wherever a quantity varies
continuously in space and/or time.
Below are representative examples across disciplines, organised by the spatial
dimension of the domain $\domain$.

\subsubsection{One-Dimensional Problems}

One-dimensional problems arise when the physical domain is a line segment
$\domain = (a,b) \subset \R$ (e.g.\ a rod, cable, or pipe cross-section).

\paragraph{Steady heat conduction in a rod.}
Given a rod of length $L$ with thermal conductivity $\kappa > 0$ and a
distributed heat source $f(x)$, the temperature $u(x)$ satisfies
\begin{equation*}
  -\kappa\,u'' = f \quad \text{in } (0,L),
  \qquad u(0) = T_0,\quad u(L) = T_L,
\end{equation*}
an \emph{elliptic} (boundary-value) problem in 1D.

\paragraph{Vibrating string (wave equation).}
The transverse displacement $u(x,t)$ of a taut string with wave speed $c$
obeys
\begin{equation*}
  u_{tt} = c^2\,u_{xx} \quad \text{in } (0,L)\times(0,T),
\end{equation*}
a \emph{hyperbolic} PDE requiring both initial displacement and velocity as
data.

\paragraph{Advection--diffusion equation.}
Transport of a scalar concentration $u(x,t)$ with velocity $b$ and
diffusivity $\varepsilon > 0$:
\begin{equation*}
  u_t + b\,u_x = \varepsilon\,u_{xx} \quad \text{in } \R\times(0,T).
\end{equation*}
This \emph{parabolic} equation models pollutant transport in rivers or
solute migration along a column.

\subsubsection{Two-Dimensional Problems}

When the domain is a planar region $\domain \subset \R^2$, the spatial
coordinates are $(x,y)$ and we encounter a rich variety of problem types.

\paragraph{Membrane deflection (Poisson equation).}
The small deflection $u(x,y)$ of a pre-stressed elastic membrane under a
transverse load $f$ satisfies
\begin{equation*}
  -\Delta u = f \quad \text{in } \domain \subset \R^2,
  \qquad u = 0 \quad \text{on } \boundary.
\end{equation*}
The same equation governs electrostatic potential in a grounded 2D conductor
and steady-state groundwater pressure.

\paragraph{Time-dependent heat conduction.}
Transient temperature in a thin plate:
\begin{equation*}
  \rho c_p\,u_t - \nabla\cdot(\kappa\,\nabla u) = f
  \quad \text{in } \domain\times(0,T),
\end{equation*}
a \emph{parabolic} PDE where $\rho$ is density and $c_p$ is specific heat.

\paragraph{Linear elasticity (plane stress/strain).}
Displacement field $\mathbf{u}(x,y) \in \R^2$ in a loaded structural
component:
\begin{equation*}
  -\nabla\cdot\boldsymbol{\sigma}(\mathbf{u}) = \mathbf{f}
  \quad \text{in } \domain,
\end{equation*}
where $\boldsymbol{\sigma}$ is the Cauchy stress tensor related to strain
by Hooke's law.
This is an \emph{elliptic system} of two coupled PDEs.

\paragraph{Shallow-water equations.}
Depth-averaged flow in coastal regions is described by the 2D
shallow-water system (a first-order \emph{hyperbolic} system) governing
water height and depth-averaged velocity.

\subsubsection{Three-Dimensional Problems}

Full 3D domains $\domain \subset \R^3$ with coordinates $(x,y,z)$ are
required when out-of-plane or depth variation cannot be neglected.

\paragraph{3D Poisson / Laplace equation.}
The electric potential $\phi$ in a dielectric medium satisfies
\begin{equation*}
  -\nabla\cdot(\varepsilon\,\nabla\phi) = \rho_e
  \quad \text{in } \domain \subset \R^3,
\end{equation*}
where $\varepsilon$ is permittivity and $\rho_e$ is charge density.
Setting $\rho_e = 0$ gives the Laplace equation, governing steady
diffusion in the absence of sources.

\paragraph{3D linear elasticity.}
Structural analysis of a 3D component under mechanical loading:
\begin{equation*}
  -\nabla\cdot\boldsymbol{\sigma}(\mathbf{u}) = \mathbf{f}
  \quad \text{in } \domain \subset \R^3,
\end{equation*}
an \emph{elliptic system} of three coupled PDEs for the displacement
components $(u_1, u_2, u_3)$.

\paragraph{Incompressible Navier--Stokes equations.}
Viscous fluid flow in 3D:
\begin{align*}
  \rho\bigl(\mathbf{u}_t + (\mathbf{u}\cdot\nabla)\mathbf{u}\bigr)
    &= -\nabla p + \mu\,\Delta\mathbf{u} + \mathbf{f}, \\
  \nabla\cdot\mathbf{u} &= 0,
\end{align*}
where $\mathbf{u}$ is velocity, $p$ pressure, $\rho$ density, and $\mu$
dynamic viscosity.
This coupled \emph{parabolic--elliptic} system governs aerodynamics,
cardiovascular blood flow, and atmospheric dynamics.

\paragraph{Maxwell's equations.}
Electromagnetic wave propagation in 3D:
\begin{align*}
  \nabla\times\mathbf{H} &= \mathbf{J} + \partial_t(\varepsilon\mathbf{E}),\\
  \nabla\times\mathbf{E} &= -\partial_t(\mu_0\mathbf{H}),
\end{align*}
a first-order \emph{hyperbolic system} for the electric field $\mathbf{E}$
and magnetic field $\mathbf{H}$, with applications in photonics, antenna
design, and MRI.

% ----------------------------------------------------------------
\subsection{What Is a PDE Mathematically?}
\label{subsec:pde-definition}

\subsubsection{Definition}

\begin{definition}[Partial Differential Equation]
A \textbf{partial differential equation} (PDE) is an equation relating an
unknown function $u \colon \domain \to \R$ (or $\R^m$) to its partial
derivatives.
For $u$ depending on independent variables
$\mathbf{x} = (x_1, \dots, x_d)$ (and possibly time $t$), a PDE of
\emph{order $k$} has the general form
\begin{equation}
  \label{eq:pde-general}
  F\!\left(\mathbf{x},\, u,\, \frac{\partial u}{\partial x_1},\dots,
    \frac{\partial^k u}{\partial x_i^{k_i}\cdots}\right) = 0,
\end{equation}
where $F$ is a given function and the highest-order derivative appearing in
$F$ determines the \emph{order} of the PDE.
\end{definition}

\subsubsection{Notation}

Throughout this document we adopt the following notational conventions.

\begin{itemize}
  \item $\domain \subset \R^d$ denotes an open, bounded domain with
    Lipschitz boundary $\boundary = \partial\domain$.
    Spatial coordinates are $\mathbf{x} = (x_1,\ldots,x_d)$, often
    written $(x)$ in 1D, $(x,y)$ in 2D, and $(x,y,z)$ in 3D.
  \item $t \in (0,T]$ is time (for time-dependent problems).
  \item Partial derivatives are written $\partial_t u \equiv u_t$,
    $\partial_{x_i} u \equiv u_{x_i}$, etc.
  \item The \textbf{gradient} of a scalar $u$ is
    $\nabla u = (\partial_{x_1}u, \ldots, \partial_{x_d}u)^T \in \R^d$.
  \item The \textbf{divergence} of a vector field
    $\mathbf{q} = (q_1,\ldots,q_d)^T$ is
    $\nabla\cdot\mathbf{q} = \sum_{i=1}^d \partial_{x_i} q_i$.
  \item The \textbf{Laplacian} of a scalar $u$ is
    $\Delta u = \nabla\cdot\nabla u = \sum_{i=1}^d \partial_{x_i}^2 u$.
  \item \textbf{Multi-index notation:} a multi-index $\alpha =
    (\alpha_1,\ldots,\alpha_d) \in \mathbb{N}_0^d$ with
    $|\alpha| = \sum_i \alpha_i$ denotes
    $D^\alpha u = \dfrac{\partial^{|\alpha|} u}{\partial x_1^{\alpha_1}
    \cdots \partial x_d^{\alpha_d}}$.
\end{itemize}

\subsubsection{Classification of Second-Order PDEs}

A second-order linear PDE in two variables $(x_1, x_2)$ with constant
coefficients has the form
$A\,u_{x_1 x_1} + B\,u_{x_1 x_2} + C\,u_{x_2 x_2} + \text{lower order} = 0$.
Its \emph{type} is determined by the discriminant $\mathcal{D} = B^2 - 4AC$:

\begin{center}
\begin{tabular}{lll}
\toprule
\textbf{Type} & \textbf{Condition} & \textbf{Prototype} \\
\midrule
Elliptic   & $B^2 - 4AC < 0$ & Laplace / Poisson equation \\
Parabolic  & $B^2 - 4AC = 0$ & Heat (diffusion) equation \\
Hyperbolic & $B^2 - 4AC > 0$ & Wave equation \\
\bottomrule
\end{tabular}
\end{center}

Each type requires different boundary and/or initial data to be
\emph{well-posed} (existence, uniqueness, and continuous dependence on
data in the sense of Hadamard).

\subsubsection{Boundary Conditions}
\label{subsubsec:bcs}

PDEs on bounded domains require \emph{boundary conditions} (BCs) on
$\boundary$ to select a unique solution.
Let $\{\Gamma_D, \Gamma_N, \Gamma_R\}$ be a non-overlapping partition of
$\boundary$ ($\Gamma_D \cup \Gamma_N \cup \Gamma_R = \boundary$).

\begin{description}
  \item[Dirichlet (essential) BC.]
    The solution value is prescribed on $\Gamma_D$:
    \begin{equation*}
      u = g_D \quad \text{on } \Gamma_D.
    \end{equation*}
    Examples: fixed temperature, clamped displacement, grounded potential.

  \item[Neumann (natural) BC.]
    The normal flux (outward derivative) is prescribed on $\Gamma_N$:
    \begin{equation*}
      \kappa\,\frac{\partial u}{\partial n}
        \equiv \kappa\,\nabla u\cdot\mathbf{n} = g_N
        \quad \text{on } \Gamma_N,
    \end{equation*}
    where $\mathbf{n}$ is the outward unit normal to $\boundary$.
    Examples: insulated boundary ($g_N = 0$), prescribed heat flux.
    The special case $g_N = 0$ is also called a \emph{no-flux} or
    \emph{free} boundary condition.

  \item[Robin (mixed) BC.]
    A linear combination of value and normal flux is prescribed on
    $\Gamma_R$:
    \begin{equation*}
      \kappa\,\frac{\partial u}{\partial n} + \alpha\,u = \beta
      \quad \text{on } \Gamma_R,
    \end{equation*}
    where $\alpha > 0$ (required to ensure the bilinear form remains
    coercive and the problem is well-posed).
    Examples: convective heat transfer (Newton cooling), impedance BC in
    acoustics.
\end{description}

\subsubsection{Initial Conditions}

For time-dependent problems ($u = u(\mathbf{x},t)$), the state of the
system at time $t=0$ must be specified.
For a first-order problem in time:
\begin{equation*}
  u(\mathbf{x}, 0) = u_0(\mathbf{x}) \quad \text{in } \domain.
\end{equation*}
Second-order hyperbolic problems (e.g.\ the wave equation) additionally
require the initial velocity:
\begin{equation*}
  u_t(\mathbf{x}, 0) = v_0(\mathbf{x}) \quad \text{in } \domain.
\end{equation*}

\subsubsection{Strong Form of a General PDE Problem}

Combining the PDE with boundary and initial conditions gives a
\emph{boundary-value problem} (BVP) or
\emph{initial-boundary-value problem} (IBVP).
As a concrete example, a general parabolic diffusion problem reads:

\begin{equation}
  \label{eq:ibvp}
  \begin{cases}
    u_t - \nabla\cdot(\kappa\,\nabla u) = f
      & \text{in } \domain\times(0,T), \\[4pt]
    u = g_D
      & \text{on } \Gamma_D\times(0,T), \\[4pt]
    \kappa\,\nabla u\cdot\mathbf{n} = g_N
      & \text{on } \Gamma_N\times(0,T), \\[4pt]
    u(\cdot,0) = u_0
      & \text{in } \domain.
  \end{cases}
\end{equation}

This \emph{strong form} requires the solution $u$ to satisfy
\eqref{eq:ibvp} pointwise.
As we shall see in Section~\ref{subsec:strong-weak}, relaxing this
pointwise requirement by multiplying by a test function and integrating
leads to the \emph{weak (variational) form}, which is the starting point
for the finite element method.

% ----------------------------------------------------------------
\subsection{Strong and Weak Formulations}
\label{subsec:strong-weak}

Consider a domain $\domain \subset \R^d$ ($d = 1, 2, 3$) with boundary
$\boundary = \partial\domain$.
A PDE in \emph{strong form} prescribes the equation pointwise throughout
$\domain$.
For example, the \textbf{Poisson equation} with homogeneous Dirichlet boundary
conditions reads:
\begin{equation}
  \label{eq:poisson-strong}
  -\Lapl u = f \quad \text{in } \domain, \qquad
  u = 0 \quad \text{on } \boundary,
\end{equation}
where $f \in L^2(\domain)$ is a given source term and $u$ is the unknown.

The strong form requires $u$ to be twice continuously differentiable, a
condition that is often too restrictive for discontinuous or rough data.
The \emph{weak} (variational) form relaxes this requirement.

\subsubsection{Deriving the Weak Form}

Multiply~\eqref{eq:poisson-strong} by a \emph{test function}
$v \in V_0 = H^1_0(\domain)$ and integrate over $\domain$:
\begin{equation*}
  -\int_\domain \Lapl u\, v \dx = \int_\domain f\, v \dx.
\end{equation*}
Applying Green's first identity (integration by parts):
\begin{equation*}
  -\int_\domain \Lapl u\, v \dx
    = \int_\domain \Grad u \cdot \Grad v \dx
      - \int_\boundary \frac{\partial u}{\partial n} v \ds,
\end{equation*}
where $\mathbf{n}$ is the outward unit normal to $\boundary$.
Since $v = 0$ on $\boundary$ (by the choice of $V_0$), the boundary term
vanishes.
The \textbf{weak formulation} of~\eqref{eq:poisson-strong} then reads:

\begin{equation}
  \label{eq:poisson-weak}
  \text{Find } u \in V \text{ such that }
  a(u, v) = L(v) \quad \forall v \in V,
\end{equation}
where the \emph{bilinear form} $a$ and the \emph{linear form} $L$ are defined
as
\begin{align}
  a(u, v) &= \int_\domain \Grad u \cdot \Grad v \dx, \label{eq:bilinear}\\
  L(v)    &= \int_\domain f\, v \dx.               \label{eq:linear}
\end{align}

\begin{remark}
In FEniCSx/UFL notation, $a(u,v)$ is called the \emph{bilinear form} or
\emph{left-hand side}, and $L(v)$ is the \emph{linear form} or
\emph{right-hand side}.
\end{remark}

\subsection{Sobolev Spaces and Function Spaces}

The natural setting for variational problems is a \emph{Sobolev space}.
\begin{definition}[Sobolev Space $H^1$]
The Sobolev space $H^1(\domain)$ consists of all square-integrable functions
whose first-order weak derivatives are also square-integrable:
\begin{equation*}
  H^1(\domain) = \bigl\{ v \in L^2(\domain)
    \mid \nabla v \in [L^2(\domain)]^d \bigr\},
\end{equation*}
equipped with the norm
$\norm{v}_{H^1}^2 = \norm{v}_{L^2}^2 + \norm{\nabla v}_{L^2}^2$.
\end{definition}

\begin{definition}[Homogeneous Dirichlet Space $H^1_0$]
$H^1_0(\domain)$ is the subspace of $H^1(\domain)$ whose elements satisfy
$v = 0$ on $\boundary$ in the trace sense.
\end{definition}

In FEniCSx, function spaces are created using the \texttt{functionspace}
function (or \texttt{FunctionSpace} in legacy FEniCS) and are characterised
by the mesh, element family, and polynomial degree.

\subsection{The Galerkin Finite Element Method}

The Galerkin method converts the infinite-dimensional problem~\eqref{eq:poisson-weak}
into a finite-dimensional approximation by replacing $V$ with a finite-dimensional
subspace $V_h \subset V$.

\subsubsection{Mesh and Triangulation}

Let $\mathcal{T}_h$ be a \emph{triangulation} (mesh) of $\domain$ into
non-overlapping elements (triangles in 2D, tetrahedra in 3D, intervals in 1D).
The subscript $h$ denotes the \emph{mesh size} (characteristic element diameter).

\subsubsection{Finite Element Basis}

On each element $K \in \mathcal{T}_h$, a polynomial space is defined.
For the standard continuous Lagrange element of degree $p$, the local space is
$\mathcal{P}_p(K)$ -- the space of polynomials of total degree at most $p$.
The global finite element space is
\begin{equation*}
  V_h = \bigl\{ v_h \in C^0(\bar\domain)
    \mid v_h|_K \in \mathcal{P}_p(K)\ \forall K \in \mathcal{T}_h \bigr\},
\end{equation*}
where continuity across element boundaries is enforced by shared degrees of
freedom (DOFs).
The homogeneous Dirichlet subspace is
$V_{h,0} = V_h \cap H^1_0(\domain)$.

\subsubsection{Discrete Variational Problem}

Replacing $V$ by $V_h$ in~\eqref{eq:poisson-weak} gives the
\emph{discrete variational problem}:
\begin{equation}
  \label{eq:discrete}
  \text{Find } u_h \in V_h \text{ such that }
  a(u_h, v_h) = L(v_h) \quad \forall v_h \in V_h.
\end{equation}

Writing $u_h = \sum_{j=1}^N U_j \phi_j$ in terms of the nodal basis
$\{\phi_j\}_{j=1}^N$ of $V_h$, problem~\eqref{eq:discrete} becomes the
linear system
\begin{equation}
  \label{eq:linsys}
  \mathbf{A}\,\mathbf{U} = \mathbf{b},
\end{equation}
where
\begin{align*}
  A_{ij} &= a(\phi_j, \phi_i) = \int_\domain \Grad\phi_j \cdot \Grad\phi_i \dx, \\
  b_i    &= L(\phi_i)        = \int_\domain f\,\phi_i \dx.
\end{align*}
The matrix $\mathbf{A}$ is the \emph{stiffness matrix} and $\mathbf{b}$ is
the \emph{load vector}.

\subsection{Boundary Conditions}

\subsubsection{Dirichlet Boundary Conditions}

A \textbf{Dirichlet} (essential) boundary condition enforces the value of
the unknown on a portion $\Gamma_D \subseteq \boundary$:
\begin{equation*}
  u = g_D \quad \text{on } \Gamma_D.
\end{equation*}
In the discrete problem, Dirichlet conditions are imposed by modifying the
linear system~\eqref{eq:linsys}.
In FEniCSx, this is done via \texttt{dirichletbc} objects passed to the
\texttt{solve} function or the \texttt{LinearProblem} class.

\subsubsection{Neumann Boundary Conditions}

A \textbf{Neumann} (natural) boundary condition prescribes the normal flux
on $\Gamma_N \subseteq \boundary$:
\begin{equation*}
  \frac{\partial u}{\partial n} = g_N \quad \text{on } \Gamma_N.
\end{equation*}
Neumann conditions appear naturally in the weak form as a boundary integral
added to the linear form:
\begin{equation*}
  L(v) = \int_\domain f\,v \dx + \int_{\Gamma_N} g_N\,v \ds.
\end{equation*}

\subsubsection{Robin (Mixed) Boundary Conditions}

A \textbf{Robin} condition combines Dirichlet and Neumann conditions:
\begin{equation*}
  \frac{\partial u}{\partial n} + \alpha u = \beta
  \quad \text{on } \Gamma_R.
\end{equation*}
This introduces both a boundary bilinear term $\alpha\int_{\Gamma_R} u v \ds$
and a boundary linear term $\int_{\Gamma_R}\beta v \ds$ into the weak form.

\subsection{A Priori Error Estimates}

A fundamental result for conforming Galerkin methods is \emph{C\'{e}a's lemma},
which states that the discrete solution is a quasi-best approximation in $V_h$:
\begin{theorem}[C\'{e}a's Lemma]
Let $u$ be the solution to the continuous problem~\eqref{eq:poisson-weak}
and $u_h$ the solution to the discrete problem~\eqref{eq:discrete}.
If $a(\cdot,\cdot)$ is continuous (bounded) and coercive on $V$ with
constants $M$ and $\alpha > 0$, then
\begin{equation*}
  \norm{u - u_h}_{H^1} \leq \frac{M}{\alpha}\,
  \inf_{v_h \in V_h} \norm{u - v_h}_{H^1}.
\end{equation*}
\end{theorem}

Combined with standard interpolation estimates, for a regular mesh and
$u \in H^{p+1}(\domain)$:
\begin{equation*}
  \norm{u - u_h}_{H^1} \leq C\, h^p\, \norm{u}_{H^{p+1}},
\end{equation*}
confirming that the error decreases at rate $h^p$ as the mesh is refined.

% ================================================================
\section{Implementation Overview}
\label{sec:implementation}
% ================================================================

This section maps the mathematical ingredients from Section~\ref{sec:math}
to concrete FEniCSx constructs.

\subsection{The FEniCSx Workflow}

A typical FEniCSx workflow follows these steps:

\begin{enumerate}
  \item \textbf{Create or import a mesh.}
  \item \textbf{Define a function space} on the mesh.
  \item \textbf{Define trial and test functions.}
  \item \textbf{Specify boundary conditions.}
  \item \textbf{Write the variational problem} using UFL.
  \item \textbf{Assemble and solve} the linear system.
  \item \textbf{Post-process and visualise} the solution.
\end{enumerate}

\subsection{Meshes in DOLFINx}

FEniCSx provides built-in mesh generators for simple domains and supports
import from external tools (Gmsh, XDMF/HDF5 files).

\begin{lstlisting}[style=python, caption={Creating a unit-square mesh in FEniCSx.}]
import dolfinx.mesh
from mpi4py import MPI
import numpy as np

# Create a unit square mesh with 32x32 quadrilateral cells
domain = dolfinx.mesh.create_unit_square(
    MPI.COMM_WORLD, 32, 32,
    cell_type=dolfinx.mesh.CellType.triangle
)
\end{lstlisting}

Key mesh attributes available in DOLFINx:
\begin{itemize}
  \item \texttt{domain.topology} -- connectivity information (vertices, edges, faces, cells).
  \item \texttt{domain.geometry} -- coordinate data.
  \item \texttt{domain.comm} -- MPI communicator (enables parallel computation).
\end{itemize}

\subsection{Function Spaces}

A \emph{finite element function space} $V_h$ is created from a mesh and an
element specification.

\begin{lstlisting}[style=python, caption={Creating a Lagrange function space.}]
from dolfinx.fem import functionspace
import basix.ufl

# Continuous Lagrange elements of degree 1 (piecewise linear)
V = functionspace(domain, ("Lagrange", 1))

# Degree-2 Lagrange for higher accuracy
V2 = functionspace(domain, ("Lagrange", 2))
\end{lstlisting}

The string \texttt{"Lagrange"} specifies the \emph{element family};
the integer is the polynomial degree $p$.
Other supported families include \texttt{"DG"} (discontinuous Galerkin),
\texttt{"N1curl"} (N\'{e}d\'{e}lec curl-conforming elements),
and \texttt{"RT"} (Raviart--Thomas divergence-conforming elements).

\subsection{Trial and Test Functions}

In UFL, \emph{trial functions} $u$ represent the unknown solution and
\emph{test functions} $v$ are used to form the variational problem.

\begin{lstlisting}[style=python, caption={Defining trial and test functions.}]
import ufl

u = ufl.TrialFunction(V)   # trial function (unknown)
v = ufl.TestFunction(V)    # test function (weight)
\end{lstlisting}

After solving, the numerical solution is stored in a \texttt{Function} object:
\begin{lstlisting}[style=python, caption={Creating a Function to hold the solution.}]
from dolfinx.fem import Function

uh = Function(V)  # will store the discrete solution u_h
\end{lstlisting}

\subsection{Variational Forms in UFL}

UFL provides a high-level syntax for expressing bilinear and linear forms
that mirrors mathematical notation.

\begin{lstlisting}[style=python, caption={Expressing the Poisson variational form in UFL.}]
import ufl
from dolfinx.fem import Constant

f = Constant(domain, 1.0)  # right-hand side f = 1

# Bilinear form a(u, v)
a = ufl.inner(ufl.grad(u), ufl.grad(v)) * ufl.dx

# Linear form L(v)
L = f * v * ufl.dx
\end{lstlisting}

\begin{table}[ht]
  \centering
  \caption{Common UFL operators and their mathematical counterparts.}
  \label{tab:ufl}
  \begin{tabular}{lll}
    \toprule
    \textbf{UFL}            & \textbf{Math}               & \textbf{Description}         \\
    \midrule
    \texttt{ufl.grad(u)}    & $\nabla u$                  & Gradient                     \\
    \texttt{ufl.div(u)}     & $\nabla \cdot \mathbf{u}$   & Divergence                   \\
    \texttt{ufl.curl(u)}    & $\nabla \times \mathbf{u}$  & Curl                         \\
    \texttt{ufl.inner(a,b)} & $a \cdot b$                 & Inner product                \\
    \texttt{ufl.dot(a,b)}   & $a \cdot b$                 & Dot product                  \\
    \texttt{ufl.dx}         & $\dx$                       & Volume integration measure   \\
    \texttt{ufl.ds}         & $\ds$                       & Exterior facet measure       \\
    \texttt{ufl.dS}         & $\dS$                       & Interior facet measure       \\
    \bottomrule
  \end{tabular}
\end{table}

\subsection{Boundary Conditions in FEniCSx}

\begin{lstlisting}[style=python, caption={Applying homogeneous Dirichlet boundary conditions.}]
import dolfinx.fem as fem
import numpy as np

# Locate all boundary facets
domain.topology.create_connectivity(
    domain.topology.dim - 1, domain.topology.dim
)
boundary_facets = dolfinx.mesh.exterior_facet_indices(domain.topology)

# Find DOFs on the boundary
boundary_dofs = fem.locate_dofs_topological(V, domain.topology.dim - 1,
                                            boundary_facets)

# Create homogeneous Dirichlet condition u = 0
bc = fem.dirichletbc(
    fem.Constant(domain, 0.0),
    boundary_dofs, V
)
\end{lstlisting}

For non-homogeneous conditions, pass a \texttt{Function} or a scalar/vector
value in place of \texttt{Constant(domain, 0.0)}.

\subsection{Assembly and Solving}

FEniCSx assembles the stiffness matrix and load vector from the UFL forms
and solves the resulting linear system.
The high-level \texttt{LinearProblem} interface handles assembly and solving
in one call:

\begin{lstlisting}[style=python, caption={Solving with LinearProblem.}]
from dolfinx.fem.petsc import LinearProblem

problem = LinearProblem(a, L, bcs=[bc],
                        petsc_options={"ksp_type": "preonly",
                                       "pc_type": "lu"})
uh = problem.solve()
\end{lstlisting}

For more control (e.g.\ custom preconditioners, iterative solvers), the
forms can be assembled separately:

\begin{lstlisting}[style=python, caption={Manual assembly using \texttt{assemble\_matrix} and \texttt{assemble\_vector}.}]
import dolfinx.fem.petsc as petsc_fem
from petsc4py import PETSc

A = petsc_fem.assemble_matrix(fem.form(a), bcs=[bc])
A.assemble()

b = petsc_fem.assemble_vector(fem.form(L))
petsc_fem.apply_lifting(b, [fem.form(a)], [[bc]])
b.ghostUpdate(addv=PETSc.InsertMode.ADD,
              mode=PETSc.ScatterMode.REVERSE)
petsc_fem.set_bc(b, [bc])
\end{lstlisting}

\subsection{Visualisation}

FEniCSx integrates with \textbf{ParaView} via XDMF/VTK output and supports
in-notebook rendering with \textbf{pyvista}.

\begin{lstlisting}[style=python, caption={Writing solution to an XDMF file for ParaView.}]
from dolfinx.io import XDMFFile

with XDMFFile(domain.comm, "poisson.xdmf", "w") as file:
    file.write_mesh(domain)
    file.write_function(uh)
\end{lstlisting}

% ================================================================
\section{Example Applications}
\label{sec:examples}
% ================================================================

\subsection{Example 1: Poisson Equation on the Unit Square}
\label{sec:ex-poisson}

We solve the Poisson equation
\begin{equation}
  -\Lapl u = f \quad \text{in } \domain = [0,1]^2,
  \qquad u = 0 \quad \text{on } \boundary,
\end{equation}
with the manufactured solution $u_\text{exact}(x,y) = \sin(\pi x)\sin(\pi y)$,
which requires $f = 2\pi^2 \sin(\pi x)\sin(\pi y)$.

\begin{lstlisting}[style=python, caption={Complete FEniCSx script for the Poisson equation.}]
import numpy as np
from mpi4py import MPI
import dolfinx
import dolfinx.mesh
from dolfinx.fem import (functionspace, Function, Constant,
                          locate_dofs_topological, dirichletbc, form)
from dolfinx.fem.petsc import LinearProblem
import ufl

# 1. Mesh
domain = dolfinx.mesh.create_unit_square(
    MPI.COMM_WORLD, 32, 32,
    cell_type=dolfinx.mesh.CellType.triangle
)

# 2. Function space
V = functionspace(domain, ("Lagrange", 1))

# 3. Trial and test functions
u = ufl.TrialFunction(V)
v = ufl.TestFunction(V)

# 4. Source term
x = ufl.SpatialCoordinate(domain)
f = 2.0 * ufl.pi**2 * ufl.sin(ufl.pi * x[0]) * ufl.sin(ufl.pi * x[1])

# 5. Variational forms
a = ufl.inner(ufl.grad(u), ufl.grad(v)) * ufl.dx
L = f * v * ufl.dx

# 6. Dirichlet boundary condition
domain.topology.create_connectivity(
    domain.topology.dim - 1, domain.topology.dim
)
boundary_facets = dolfinx.mesh.exterior_facet_indices(domain.topology)
boundary_dofs = locate_dofs_topological(
    V, domain.topology.dim - 1, boundary_facets
)
bc = dirichletbc(Constant(domain, 0.0), boundary_dofs, V)

# 7. Solve
problem = LinearProblem(a, L, bcs=[bc],
                        petsc_options={"ksp_type": "preonly",
                                       "pc_type": "lu"})
uh = problem.solve()

# 8. Compute L2 error
u_exact = ufl.sin(ufl.pi * x[0]) * ufl.sin(ufl.pi * x[1])
error_form = form((uh - u_exact)**2 * ufl.dx)
import dolfinx.fem
error_L2 = np.sqrt(
    domain.comm.allreduce(
        dolfinx.fem.assemble_scalar(error_form), op=MPI.SUM
    )
)
if domain.comm.rank == 0:
    print(f"L2 error: {error_L2:.3e}")
\end{lstlisting}

\subsection{Example 2: Transient Heat Equation}
\label{sec:ex-heat}

The heat equation models time-dependent diffusion:
\begin{equation}
  \frac{\partial u}{\partial t} - \Lapl u = f
  \quad \text{in } \domain \times (0, T],
  \qquad u = 0 \text{ on } \boundary,
  \qquad u(\cdot, 0) = u_0 \text{ in } \domain.
\end{equation}

Applying the \textbf{backward Euler} time-stepping scheme with step $\Delta t$:
\begin{equation*}
  u^{n+1} - u^n = \Delta t\,\bigl(\Lapl u^{n+1} + f^{n+1}\bigr),
\end{equation*}
the weak form at each time step becomes:
\begin{equation}
  \underbrace{\int_\domain u^{n+1} v \dx
    + \Delta t \int_\domain \Grad u^{n+1} \cdot \Grad v \dx}_{a(u^{n+1},v)}
  = \underbrace{\int_\domain u^n v \dx
    + \Delta t \int_\domain f^{n+1} v \dx}_{L(v)}.
\end{equation}

\begin{lstlisting}[style=python, caption={Time-stepping loop for the heat equation (FEniCSx).}]
# (Continuing from Example 1 setup -- mesh and V already created)
import numpy as np
from dolfinx.fem import Function, Constant
from dolfinx.fem.petsc import LinearProblem
import ufl

dt = 0.01
T  = 0.5

u_n = Function(V)    # solution at previous time step
u_n.interpolate(lambda x: np.sin(np.pi * x[0]) * np.sin(np.pi * x[1]))

u  = ufl.TrialFunction(V)
v  = ufl.TestFunction(V)
f  = Constant(domain, 0.0)

a = (u * v + dt * ufl.inner(ufl.grad(u), ufl.grad(v))) * ufl.dx
L = (u_n * v + dt * f * v) * ufl.dx

t = 0.0
while t < T:
    t += dt
    problem = LinearProblem(a, L, bcs=[bc],
                            petsc_options={"ksp_type": "preonly",
                                           "pc_type": "lu"})
    uh = problem.solve()
    u_n.x.array[:] = uh.x.array
\end{lstlisting}

\subsection{Example 3: Linear Elasticity}
\label{sec:ex-elasticity}

Linear elasticity governs the displacement $\mathbf{u}$ of an elastic body:
\begin{equation*}
  -\Div\boldsymbol{\sigma}(\mathbf{u}) = \mathbf{f}
  \quad \text{in } \domain,
\end{equation*}
where the \emph{stress tensor} $\boldsymbol{\sigma}$ is related to the
\emph{strain tensor} $\boldsymbol{\varepsilon}$ through
\emph{Hooke's law} (isotropic material):
\begin{align*}
  \boldsymbol{\sigma}(\mathbf{u})
    &= \lambda (\Div\mathbf{u})\mathbf{I}
      + 2\mu\,\boldsymbol{\varepsilon}(\mathbf{u}),\\
  \boldsymbol{\varepsilon}(\mathbf{u})
    &= \tfrac{1}{2}\bigl(\nabla\mathbf{u} + (\nabla\mathbf{u})^T\bigr),
\end{align*}
with Lam\'{e} parameters $\lambda, \mu > 0$.

The variational form for vector-valued displacement $\mathbf{u}$ is:
\begin{equation}
  a(\mathbf{u},\mathbf{v})
    = \int_\domain \boldsymbol{\sigma}(\mathbf{u}) : \boldsymbol{\varepsilon}(\mathbf{v}) \dx
    = \int_\domain \mathbf{f} \cdot \mathbf{v} \dx.
\end{equation}

In FEniCSx, vector-valued problems use a \emph{vector element}:
\begin{lstlisting}[style=python, caption={Setting up a vector function space for elasticity.}]
# Vector Lagrange space of degree 1 in 2D
V_vec = functionspace(domain, ("Lagrange", 1, (domain.geometry.dim,)))

lam = Constant(domain, 1.0)
mu  = Constant(domain, 1.0)

def epsilon(u):
    return ufl.sym(ufl.grad(u))

def sigma(u):
    return lam * ufl.div(u) * ufl.Identity(len(u)) + 2 * mu * epsilon(u)

u = ufl.TrialFunction(V_vec)
v = ufl.TestFunction(V_vec)
f = ufl.as_vector([0.0, -1.0])

a = ufl.inner(sigma(u), epsilon(v)) * ufl.dx
L = ufl.inner(f, v) * ufl.dx
\end{lstlisting}

% ================================================================
\section*{Conclusion}
\addcontentsline{toc}{section}{Conclusion}
% ================================================================

This document has introduced the theoretical foundations of the Finite Element
Method as implemented in FEniCSx.
The key concepts are:
\begin{itemize}
  \item Derivation of the \emph{weak (variational) form} by multiplying the
    strong PDE by test functions and integrating by parts.
  \item The \emph{Galerkin discretisation} that restricts the problem to a
    finite-dimensional subspace $V_h$ defined by a mesh and a polynomial degree.
  \item FEniCSx abstractions: \texttt{mesh}, \texttt{FunctionSpace},
    \texttt{TrialFunction}/\texttt{TestFunction}, \texttt{dirichletbc},
    and UFL forms assembled by \texttt{LinearProblem}.
  \item Worked examples: Poisson equation, heat equation (transient),
    and linear elasticity.
\end{itemize}

For further reading, the official FEniCSx tutorial~\cite{Dokken2023} and
the FEniCS book~\cite{LoggMardalWells2012} are excellent starting points.

% ================================================================
\begin{thebibliography}{9}
% ================================================================

\bibitem{LoggMardalWells2012}
  A.~Logg, K.-A.~Mardal, G.~N.~Wells (eds.),
  \textit{Automated Solution of Differential Equations by the Finite Element Method},
  Springer, 2012.
  \href{https://doi.org/10.1007/978-3-642-23099-8}{doi:10.1007/978-3-642-23099-8}

\bibitem{Scroggs2022}
  M.~W.~Scroggs, J.~S.~Dokken, C.~N.~Richardson, G.~N.~Wells,
  Construction of Arbitrary Order Finite Element Degree-of-Freedom Maps on
  Polygonal and Polyhedral Cell Meshes,
  \textit{ACM Transactions on Mathematical Software}, 48(2), 2022.
  \href{https://doi.org/10.1145/3524456}{doi:10.1145/3524456}

\bibitem{Dokken2023}
  J.~S.~Dokken,
  \textit{The FEniCSx Tutorial},
  Simula Research Laboratory, 2023.
  \url{https://jsdokken.com/dolfinx-tutorial/}

\bibitem{Brenner2008}
  S.~C.~Brenner, L.~R.~Scott,
  \textit{The Mathematical Theory of Finite Element Methods},
  3rd ed., Springer, 2008.

\bibitem{Ern2004}
  A.~Ern, J.-L.~Guermond,
  \textit{Theory and Practice of Finite Elements},
  Springer, 2004.

\end{thebibliography}

\end{document}
